\subsection{JSON-Struktur für den Graph}

Damit man den STEMgraph später als Graph erkennen kann, brauche ich ein JSON, das
genau zu meinem Datenmodell passt. Ich habe mir das so aufgebaut:

\begin{verbatim}
{
  "nodes": [
    { "id": "7392771207492170471730773974971",
      "teaches": "Linux CLI Grundlagen", "keywords": ["linux", "cli"] },
    { "id": "27498625678403678296478352674894",
      "teaches": "Dateien und Ordner verwalten", "keywords": ["linux", "files", "mkdir"] },
    { "id": "23684937628947256748937849374783",
      "teaches": "Makefile Grundlagen", "keywords": ["makefile", "buildsystem"] }
  ],
  "edges": [
    { "source": "7392771207492170471730773974971",
      "target": "27498625678403678296478352674894" },
    { "source": "27498625678403678296478352674894",
      "target": "23684937628947256748937849374783" }
  ]
}
\end{verbatim}
Die \texttt{nodes} sind die Challenges mit meinen Feldern \texttt{id}, \texttt{teaches} und \texttt{keywords}.
\\
Die \texttt{edges} sind die Abhängigkeiten und bilden \texttt{depends\_on} ab, indem eine
Challenge als \texttt{source} und die Challenge, die darauf aufbaut, als \texttt{target}
eingetragen wird.